\chapter{Workspace Architecture}

\section{Repository Map}

The repository separates mathematical claims, computational evidence, source
code, paper text, and private notes. This separation is part of the research
method: a computation may suggest a claim, but it does not by itself prove it.

\section{Theory And Claims}

\path{theory/} stores the evolving mathematical development. The directory
\path{theory/claims/} is the formal claim registry, and every theorem-like
statement should eventually be represented there with status, assumptions,
dependencies, proof sketch, references, and verification notes.

The current claim statuses are deliberately conservative. \path{CLAIM-0001}
and \path{CLAIM-0002} are draft claims, while \path{CLAIM-0003} through
\path{CLAIM-0005} are conjectural.

\section{Code, Tests, And Experiments}

\path{src/jordan_qh/} contains exact Python scaffolding for small examples,
linear algebra, identities, derivations, modules, omega candidates,
indecomposables, and toy Quillen checks. The matching tests live in
\path{tests/}.

\path{experiments/} records planned or executed computations. The experiment
numbers are stable workspace identifiers, not a logical dependency order.

\section{Paper, Formalization, And Private Notes}

\path{paper/} is the public-facing LaTeX skeleton and should only receive
statements that are aligned with checked or clearly labelled claim statuses.
\path{formal/} is a Lean scaffold, not part of the default verification path.

\path{note/} is a private learning area. This file may contain bilingual
explanations, partial understanding, and personal TODOs, but it should not
silently promote experiments or conjectures into proved results.
